\chapter{Architecture Overview}
\ChapterLabel{ArchitectureOverview}
\section{Chapter Responsibility}
\label{Section:ArchitectureOverviewResponsibility}
Explain the major design choices and why they address the design drivers.

\noindent\textbf{Expected result:} A defensible design direction with explicit trade-offs.

\section{How to Use This Chapter}
Explain the forces that shape the design, the major decisions, alternatives considered, and consequences. Use the 4+1 views only when they help readers answer different architectural questions.
\section{Proposed Local and OMS Boundary Decision}
\label{Section:ProposedOMSBoundary}
\ArchitectureID{ARC-002}: Proposed, not approved. Keep the local sample
record contract (\RequirementRef{REQ-IF-01}) distinct from the mission-facing
derived sensor-status contract (\RequirementRef{REQ-OMS-01}). The former
carries acquisition records; the latter depends on OMS applicability,
version, CAL profile, UCI mapping, gateway ownership, and service decisions.
Alternatives are a selected OMS boundary or another explicitly approved
mission contract; neither has been selected here. Separate contracts permit
local acquisition evolution but require an explicit mapping and end-to-end
freshness/failure allocation. \hyperref[Parameter:PAR-07]{PAR-07} in
\hyperref[Table:RequirementParameters]{the decision register} gates approval
and conformance testing. No driver, FIFO mode, latch, or counter part is
selected by this decision proposal.

\section{Decision Template and Example}
\textbf{ARC-001:} \textit{[Decision]}; \textbf{Drivers:} \textit{[quality attributes and constraints]}; \textbf{Rationale:} \textit{[why selected]}; \textbf{Consequences:} \textit{[benefits, trade-offs, and risks]}. Example: use an asynchronous message queue to protect request processing from transient downstream outages; this improves resilience but introduces eventual consistency.
